\clearpage\thispagestyle{empty}
\pdfbookmark[0]{Quellen zu Juno (9. September 1812)}{juno1812sept9-sources}
\begin{center}\textbf{Quellen zu Juno (9. September 1812)}\end{center}
\noindent Diese KI-Übertragung ins heutige Deutsch folgt dem Brief vom 9. September 1812 in den \href{https://gdz.sub.uni-goettingen.de/id/PPN236007467}{\textit{Werken}, Band~VI} von Gauß (Göttingen, 1874), S.~360--361. Der Brief erschien zuerst als Nr.~XXIV in der Zeitschrift \textit{Monatliche Correspondenz zur Beförderung der Erd- und Himmelskunde}, Band~XXVI (September 1812), S.~\href{https://api.digitale-sammlungen.de/iiif/image/v2/bsb10538615_00307/full/1600,/0/default.jpg}{297}, \href{https://api.digitale-sammlungen.de/iiif/image/v2/bsb10538615_00308/full/1600,/0/default.jpg}{298}, \href{https://api.digitale-sammlungen.de/iiif/image/v2/bsb10538615_00309/full/1600,/0/default.jpg}{299}. Die Links führen zum digitalisierten Exemplar der Bayerischen Staatsbibliothek.
\par\smallskip\noindent Der Erstdruck wiederholt Monatsangaben und Deklinationsrichtung durch Wiederholungsstriche, gibt als Rektaszensionsdifferenz für den 12. Juli $+22{,}4$ an und bezeichnet die geozentrische Breite der Opposition als nördlich. Die Werkausgabe lässt diese Wiederholungsstriche weg, gibt $+22\;4$ an und lässt die Richtung im Oppositionsergebnis offen. Bei der früheren Überschneidung der Bahnen nennen beide Ausgaben die vierte und sechste Spalte. KI-Übertragung und KI-Satz.
\clearpage
